% BHU PYQ -- Manually transcribed & error-corrected
\documentclass[11pt,a4paper]{article}

\usepackage[a4paper,top=2.5cm,bottom=2.5cm,left=2.8cm,right=2.8cm,headheight=14pt]{geometry}
\usepackage{amsmath,amssymb}
\usepackage[T1]{fontenc}
\usepackage[utf8]{inputenc}
\usepackage{lmodern}
\usepackage{microtype}
\usepackage{fancyhdr}
\usepackage{enumitem}
\usepackage{hyperref}
\usepackage{lastpage}
\usepackage{parskip}

\hypersetup{colorlinks=true,urlcolor=black,linkcolor=black,
  pdftitle={CHB-504: Physical Chemistry-III -- BHU PYQ},pdfauthor={Banaras Hindu University}}

\pagestyle{fancy}\fancyhf{}
\renewcommand{\headrulewidth}{0.4pt}\renewcommand{\footrulewidth}{0.4pt}
\fancyhead[L]{\small   \textit{Banaras Hindu University}}
\fancyhead[R]{\small   \textit{CHB-504: Physical Chemistry-III}}
\fancyfoot[C]{\small Page  \thepage\ of \pageref{LastPage}}


\newlist{parts}{enumerate}{2}
\setlist[parts,1]{label=(\alph*),leftmargin=*,itemsep=5pt,topsep=3pt}
\setlist[parts,2]{label=(\roman*),leftmargin=*,itemsep=3pt,topsep=2pt}

\newcommand{\pts}[1]{\hfill   \textit{[#1]}}


\begin{document}


\begin{center}
  {\large\bfseries BANARAS HINDU UNIVERSITY}\\[2pt]
  {\normalsize B.Sc. (Hons.) Semester V Examination, 2022-23}\\[6pt]
  \rule{0.8\linewidth}{0.5pt}\\[6pt]
  {\large\bfseries Paper No. CHB-504: Physical Chemistry-III}\\[4pt]
  {\normalsize Time: 3 Hours\hfill Maximum Marks: 70}
\end{center}
\rule{\linewidth}{0.4pt}
\small



\noindent   \textit{Instructions: Answer five questions in total, including Question~1 which is compulsory.}

\normalsize
\bigskip




\noindent   \textbf{Question 1.} Answer all parts of the following: \pts{5 $ \times$ 2 = 10}

\begin{parts}
  \item Brief about any three processes of deactivation of a photochemically activated molecule (internal conversion, intersystem crossing, fluorescence, etc.).
  \item Give two experimental evidences which validate the Debye-H\"uckel model of electrolyte solutions.
  \item Explain how the chemical potential of a reference ion changes as a consequence of ion-ion interactions.
  \item Define fugacity and describe its physical significance.
  \item Write the relationship between activity and mole fraction.
\end{parts}

\medskip\rule{\linewidth}{0.2pt}




\noindent   \textbf{Question 2.}

\begin{parts}
  \item What is an ideal-dilute solution? Explain with the help of Henry's and Raoult's laws. \pts{6}
  \item State: (i) Gibbs-Duhem equation and (ii) Duhem-Margules equation, and explain their thermodynamic significances. \pts{8}
\end{parts}

\medskip\rule{\linewidth}{0.2pt}




\noindent   \textbf{Question 3.}

\begin{parts}
  \item State the physical significance of fugacity. Describe the Lewis-Randall rule for gas mixtures. \pts{6}
  \item Describe osmometry and its applications in the determination of molar masses of macromolecules. \pts{8}
\end{parts}

\medskip\rule{\linewidth}{0.2pt}




\noindent   \textbf{Question 4.}

\begin{parts}
  \item If $I_a$ is the intensity of the light absorbed, derive an expression for the quantum yield of the photodimerization of anthracene ($\Phi$). \pts{8}
  \item Discuss the influence of pH and temperature, depicting the graph, on the kinetics of enzyme-catalyzed reactions. \pts{6}
\end{parts}

\medskip\rule{\linewidth}{0.2pt}




\noindent   \textbf{Question 5.}

\begin{parts}
  \item What do you understand by the thickness or radius of the ionic cloud that surrounds a reference ion? Show the variation of the thickness of the ionic cloud as a function of electrolyte concentration and prove that it has the dimensions of length. \pts{8}
  \item Under which conditions is the Wien effect observed? Explain, giving examples, the increase in the conductivities of electrolytic solutions with the application of high electric fields. \pts{6}
\end{parts}

\vfill
\rule{\linewidth}{0.4pt}

\begin{center}\small
  \href{https://bhu-exam.vercel.app/}{Click here to visit our website}\\[4pt]
  \href{https://me-aryan.vercel.app/}{Click here to visit our contributor}\\[4pt]
  \href{https://quashan-me.vercel.app/}{Click here to visit our contributor}
\end{center}

\end{document}
